\documentclass[11pt,a4paper,twoside,openright]{book}

% ── Γραμματοσειρές & κωδικοποίηση ────────────────────────────
\usepackage[T1]{fontenc}
\usepackage[utf8]{inputenc}
\usepackage[greek,english]{babel}   % κύρια γλώσσα: ελληνικά
\usepackage{alphabeta}             % ενεργοποιεί \alpha, \beta κ.λπ. με babel

% ── Μαθηματικά ──────────────────────────────────────────────
\usepackage{amsmath,amssymb,amsthm}
\usepackage{mathtools}

% ── Γεωμετρία σελίδας ───────────────────────────────────────
\usepackage[a4paper,inner=3cm,outer=2.5cm,top=2.5cm,bottom=2.5cm]{geometry}

% ── Ορισμός θεωρηματικών περιβαλλόντων (εργονομική αρίθμηση) ─
\newtheorem{lemma}{Λήμμα}[section]          % αρίθμηση ανά section
\newtheorem{proposition}[lemma]{Πρόταση}    % μοιράζεται μετρητή με το lemma
\newtheorem{theorem}[lemma]{Θεώρημα}        % μοιράζεται μετρητή με το lemma
\newtheorem{observation}[lemma]{Παρατήρηση} % μοιράζεται μετρητή με το lemma

% Τα «proof» περιβάλλοντα μένουν χωρίς αρίθμηση
\renewcommand{\proofname}{\textnormal{Απόδειξη}}

% ── Ενδεικτικές εντολές φυσικής ─────────────────────────────
\newcommand{\dd}[1]{\,\mathrm{d}#1}         % διαφορικό
\newcommand{\R}{\mathbb{R}}                 % πραγματικοί αριθμοί

% ═════════════════════════════════════════════════════════════
\begin{document}

\title{Μαθηματικές Μέθοδοι Φυσικής}
\author{Συγγραφέας}
\date{\today}
\maketitle

\chapter{Στοιχεία Διαφορικής Γεωμετρίας}

\section{Λείες Πολλαπλότητες --- Βασικές Έννοιες}

\begin{lemma}[Θεώρημα αντίστροφης απεικόνισης]
  Έστω $f: \R^n \to \R^n$ λεία συνάρτηση και $p\in\R^n$ με
  $\det Df(p) \neq 0$. Τότε υπάρχουν ανοικτές περιοχές $U\ni p$, $V\ni f(p)$
  ώστε η $f|_U : U\to V$ να είναι αμφιδιαφόριση.
\end{lemma}

\begin{proposition}[Θεώρημα Nash]
  Κάθε λεία πολλαπλότητα $M$ εμφυτεύεται ισομετρικά σε κάποιον
  Ευκλείδειο χώρο $\R^N$.
\end{proposition}

\begin{theorem}[Stokes]\label{thm:stokes}
  Για κάθε προσανατολισμένη $n$-διάστατη πολλαπλότητα $M$ με σύνορο
  $\partial M$ και κάθε $(n-1)$-μορφή $\omega$ συμπαγούς φορέα ισχύει
  \begin{equation}\label{eq:stokes}
    \int_M \dd{\omega} = \int_{\partial M} \omega.
  \end{equation}
\end{theorem}

\begin{observation}
  \label{obs:stokes-sign}
  Στη φυσική βιβλιογραφία το θεώρημα Stokes εμφανίζεται συχνά με αντίθετο
  πρόσημο, ανάλογα με τη σύμβαση προσανατολισμού.  Στο εξής θα υιοθετήσουμε
  τη μαθηματική σύμβαση της εξίσωσης \eqref{eq:stokes}.
\end{observation}

\end{document}